\documentclass[11pt]{article}
\usepackage[utf8]{inputenc}
\usepackage{geometry}
\geometry{a4paper, margin=1in}
\usepackage{booktabs}
\usepackage{enumitem}
\usepackage{hyperref}

\title{\textbf{Kenya Informatics Olympiad (KIO) \\ Operational Plan \& Selection Strategy}}
\author{KIO Organizing Committee}
\date{August 2026}

\begin{document}

\maketitle

\section{Introduction \& Objectives}
The Kenya Informatics Olympiad (KIO) is the national program designed to identify, train, and select secondary school students to represent Kenya in international competitive programming. 

Kenya is hosting the \textbf{Pan-African Informatics Olympiad (PAIO)} in \textbf{August 2027}. This hosting responsibility gives us a clear deadline: we must select, train, and prepare a highly competitive four-member national team by July 2027.

\subsection{Eligibility Criteria}
To ensure continuity, training longevity, and smooth transitions, we are restricting participation to students in \textbf{Forms 1, 2, and 3}.

We will exclude \textbf{Form 4} students. Because they sit their national exams (KCSE) at the end of the year and transition to university, they cannot participate in the extensive training camp or represent the country in the following year's international competitions.

\section{The Organizing Team}
Our leadership consists of two directors and four dedicated coaches, with plans to add more volunteers as the program grows.

\subsection{Directors}
\begin{itemize}
    \item \textbf{Dr. James Katende} (University of Nairobi) -- Academic oversight and university partnerships.
    \item \textbf{Martin Mungai} (CEMASTEA) -- School outreach and government coordination.
\end{itemize}

\subsection{Coaches \& Operations}
\begin{itemize}
    \item \textbf{Dalton Omondi} (Kenyatta University) -- Head Coach, curriculum and contest design.
    \item \textbf{Angelo Opiyo} (JKUAT) -- Platform engineer and training coordinator.
    \item \textbf{Bill Mbogho} (Huawei) -- Industry liaison and technical mentor.
    \item \textbf{Ahona Yoshi} (Arizona State University) -- International training coordinator.
    \item \textit{[Placeholder for Operations Support]} -- (Note: Please add the name and details of the operations team member you added previously here. We sincerely apologize for overwriting your edit on the last file save!)
\end{itemize}

\section{Standard Structure of the Selection Rounds}
To establish a clear reference point, KIO's selection pipeline is built on three distinct testing rounds, each designed to evaluate different stages of problem-solving and programming proficiency:

\begin{description}
    \item[Round 1: Logic \& Computational Thinking] \hfill \\
    An online qualifier consisting of multiple-choice questions (options A, B, C, D, E). This round requires no prior programming knowledge and focuses purely on logic, pattern recognition, and basic algorithmic thinking.
    
    \item[Round 2: Algorithmic Description \& Pseudocode] \hfill \\
    A written paper testing design logic. Students are given problems and must describe their step-by-step algorithms. Partial marks are awarded for clear pseudocode or logical descriptions, while full marks are reserved for writing actual syntax-correct code. This allows brilliant beginners to compete.
    
    \item[Round 3: Pure Programming Contest] \hfill \\
    A live, automated programming contest hosted on Codeforces. Students must submit working code that passes automated test cases under time limits. Programming proficiency is mandatory.
\end{description}

\section{Immediate Cycle: PAIO 2027 Selection}
Because the timeline between now (late August 2026) and the August 2027 competition is extremely tight, we must compress the selection timeline. 

\subsection{Timeline and Scheduling Constraints}
In the third school term, when the KCSE national exams begin in late November, Form 1, 2, and 3 students will be sent home, leaving only Form 4 students in school. This means December is the ideal window for our first qualifier:
\begin{itemize}
    \item \textbf{Timing Constraints}: The annual Kenya Mathematical Olympiad (KMO) Round 3 camp starts around \textbf{December 16th}. Since some of our prospective students will be attending the KMO camp, we must complete our qualifier before this date.
    \item \textbf{Proposed Date}: We propose \textbf{Saturday, December 5th, 2026} (or alternatively \textbf{Wednesday, December 9th, 2026}). We must avoid Saturday, December 12th, as it is Jamhuri Day (a major national holiday in Kenya).
    \item \textbf{Format \& Logistics}: The test will be run online using custom proctoring software while students are at home for the holidays. The advantage is that students will have no schoolwork burden; the challenge is ensuring they have internet access at home.
\end{itemize}

We are considering two potential operational options for this immediate cycle:

\subsection{Option A: Compressed 3-Round Track}
We execute a fast-tracked version of all three standard rounds:
\begin{enumerate}
    \item \textbf{First Paper (December 5th, 2026)}: Online Round 1 exam (Multiple-choice A--E).
    \item \textbf{Second Paper (February 2027)}: Round 2 exam (Algorithmic description + code), completed in-person/at hubs no later than \textbf{February 28th, 2027}.
    \item \textbf{Third Paper (April 2027)}: Round 3 exam (Pure Codeforces) administered during the Rising Stars Camp hosted by the African Olympiad Academy (AOA).
\end{enumerate}

\subsection{Option B: 2-Round Direct Track}
We skip the multiple-choice qualifying round entirely and begin selection with the second-stage exam:
\begin{enumerate}
    \item \textbf{First Paper (December 5th, 2026)}: Run a direct Round 2 exam (Algorithmic description + code) online with proctoring.
    \item \textbf{Second Paper (April 2027)}: Select the top performers and invite them directly to the AOA Rising Stars Camp for the final Round 3 Codeforces exam.
\end{enumerate}

\section{Subsequent Cycles (2028 and Beyond)}
For future years, we will establish a comfortable, non-rushed schedule by initiating the selection rounds in the middle of the school year. 

\subsection{Proposed 3-Round Timeline}
\begin{enumerate}
    \item \textbf{Round 1 (June/July -- Term 2)}: Online qualifier run on a designated weekend while students are in school. This allows school computer labs to be utilized under teacher supervision.
    \item \textbf{Round 2 (December)}: In-person written semifinal testing algorithmic logic and code structure, held after school terms close.
    \item \textbf{Round 3 (April Holiday)}: Intensive coding camp and final selection on Codeforces.
\end{enumerate}

\section{Funding \& Logistics Challenges}
\begin{enumerate}
    \item \textbf{Funding Uncertainty}: Our budget for physical camps next year is not yet guaranteed.
    \begin{itemize}
        \item \textit{If funded}: We will run physical camps and align our timelines with the Kenya Mathematical Olympiad (KMO) to share travel and venue costs.
        \item \textit{If unfunded}: We will host virtual intensive bootcamps, and run Round 2 via local university hubs (Nairobi, Kisumu, Mombasa) to eliminate student travel expenses.
    \end{itemize}
    \item \textbf{Bridging the Syntax Gap}: Transitioning students from logic puzzles (Round 1) to actual programming (Round 2) is difficult. The coaches will organize a free virtual coding bootcamp between January and March for all Round 1 qualifiers.
\end{enumerate}

\end{document}
